\documentclass[11pt,a4paper]{article}
\usepackage[margin=2.4cm]{geometry}
\usepackage{amsmath,amssymb,booktabs,graphicx,hyperref,natbib}
\usepackage{siunitx}
\newcommand{\Eh}{E_h}
\newcommand{\bohr}{\mathrm{bohr}}
\graphicspath{{../mouse_orbital_optimized_dft/figures/}}
\title{Orbital-Optimized DFT for Excited States in MOUSE:$\Delta$SCF, Spin-Pure ROKS, and Analytic Nuclear Response}
\author{Vincenzo Barone}
\date{Draft companion manuscript --- 2 November 2026}
\begin{document}
\maketitle

\begin{abstract}
We present a full-integral DFT implementation of state-specific excited-state
optimization in MOUSE. Fixed-occupation unrestricted $\Delta$SCF and a
common-orbital two-open-orbital ROKS construction are implemented on top of
the ground-state PySCF machinery. The ROKS energy is the multiplet combination
$E_S=2E_M-E_T$ evaluated with common spatial orbitals, giving a coherent
spin-pure singlet representation. Orbital stationarity is solved with a
residual trust-region algorithm and its coupled response is differentiated
analytically through molecular gradients and Hessians, including moving-grid
XC terms. The present draft establishes the implementation and derivative
qualification; chemical accuracy, broad transferability, and a complete
comparison with ORCA MOM/GMF and native PySCF MOM remain future work.
\end{abstract}

\section{Scope and motivation}
Response methods such as TDA and TD-DFT describe many excited states from a
ground-state reference. State-specific orbital optimization instead obtains an
excitation from the difference between stationary electronic energies and can
reuse the energy and derivative infrastructure of a ground-state program.
Levi's direct-optimization and saddle-search work motivates treating an
excited state as a stationary point, often a saddle, rather than as a fragile
SCF fixed point. The present work develops that idea as a general full-DFT
backend, separate from the semiempirical OO-CID model reported elsewhere.

\section{Theory}
For a fixed non-Aufbau occupation, unrestricted $\Delta$SCF optimizes the
spin orbitals directly and may remain spin contaminated. For a two-open-
orbital singlet, MOUSE uses common spatial orbitals and evaluates
\begin{equation}
 E_S[\mathbf C]=2E_M[\mathbf C]-E_T[\mathbf C],
\end{equation}
where mixed and triplet components share the same orbitals. The resulting
two-determinant CSF has $\langle S^2\rangle=0$ by construction, while the
DFT multiplet-sum energy remains an approximate functional expression.

\section{Implementation}
The implementation is a PySCF extension with fixed occupations, state
following, orbital continuation, and optional freeze-and-release preparation.
The final solve releases all nonredundant orbital rotations. The residual
Jacobian is treated as nonsymmetric away from stationarity; its exact adjoint
drives a small residual trust-region step. Range-separated exact integrals
are cached when memory permits, with an exact native fallback.

\subsection{Full-DFT implementation of the orbital response}
A separate reference implementation reuses molecular UHF/UKS energies,
Fock matrices and analytic density-response kernels from PySCF.\cite{sun2020pyscf}
Its unrestricted $\Delta$SCF and common-orbital ROKS solvers keep integer
occupations fixed and solve orbital stationarity with Newton--Krylov
steps, retaining negative curvature. They require no TDA calculation or
MOM occupation reassignment. The ROKS wavefunction is a coherent pure
singlet; its $2E_M-E_T$ functional energy is evaluated with common spatial
orbitals and remains an approximation. Neither an energy offset nor the
semiempirical kernel parameters enter this full-DFT implementation.

The nuclear response includes the complete noncanonical energy-weighted
density and one coupled common-orbital ROKS CP equation. Analytic moving
Becke-grid derivatives, Gaussian AO derivatives through third order and
LibXC vxc/fxc supply the XC contribution to the Hessian. Independent
checks on water and formaldehyde cover HF, LDA, PBE0 and CAM-B3LYP; the
DFT Hessian--vector errors are below $5.1\times10^{-9}$ Ha/bohr$^2$.
The SI reports the separate numerical certificates. These small-basis
checks validate implementation, rather than spectral accuracy or a
speed advantage over MOM. They do not extend the semiempirical
26-state accuracy ledger. Reusing spin-density/AO products accelerates
the grid chain rule without changing the complete Hessians by more than
$2.4\times10^{-13}$ Ha/bohr$^2$ in these checks.


\section{Qualification and current limitations}
The qualified scope is all-electron, full-integral molecular HF and LDA/GGA,
including global and range-separated hybrids, with analytic moving-grid
response. Density fitting, solvent, relativistic, meta-GGA and ECP Hessians
are outside the present qualification. The tests establish derivative
consistency and solver robustness; they do not yet establish a general
accuracy advantage over MOM, ORCA, PySCF, TDA, or other response methods.

\section{Outlook}
The next stage will compare identical roots and functionals across MOUSE,
ORCA and native PySCF, add GMF/saddle-following benchmarks, and evaluate
excited-state geometries, frequencies, CT and Rydberg states on a frozen
independent set.

\bibliographystyle{unsrtnat}
\bibliography{../mouse_orbital_optimized_dft/references}
\end{document}
